% Generated draft: requires mathematical and source review.
\documentclass[11pt]{article}
\usepackage[margin=25mm]{geometry}
\usepackage{fontspec}
\setmainfont{Latin Modern Roman}
\usepackage{amsmath,amssymb,mathtools,unicode-math}
\setmathfont{Latin Modern Math}
\usepackage{enumitem,needspace,xurl,hyperref}
\setlength{\parindent}{0pt}
\setlength{\parskip}{5pt}
\setlength{\emergencystretch}{3em}
\newcommand{\N}{\mathbb{N}}
\newcommand{\Z}{\mathbb{Z}}
\newcommand{\Q}{\mathbb{Q}}
\newcommand{\R}{\mathbb{R}}
\newcommand{\C}{\mathbb{C}}
\newcommand{\F}{\mathbb{F}}
\newcommand{\A}{\mathbb{A}}
\newcommand{\PP}{\mathbb{P}}
\newcommand{\E}{\mathbb{E}}
\newcommand{\eps}{\varepsilon}
\newcommand{\dd}{\,\mathrm{d}}
\begin{document}
\section*{UnsolvedMath ID: 20002012}
\textbf{Source problem code:} catalogue entry
\subsection{Definitions and mathematical statement}
The exact canonical Conjecture 2 says: ``Any \(S\) as above may be written \(\mathrm{const}\cdot Q+L+V\), where \(V\) is an exact divergence.'' The antecedent defining \(S\), the dimension, and the meanings and normalization of \(Q,L\) are omitted. The supplied aid interprets this as a decomposition of natural top-degree forms with conformally invariant integrals into a constant multiple of a \(Q\)-curvature density, a local conformal invariant, and an exact divergence. This interpretation is explicit rather than a recovery of all missing source hypotheses. In a closed oriented setting, an exact top-degree form \(V=\dd\eta\) integrates to zero; a divergence density \((\operatorname{div}_gX)\dd V_g\) has this form.

The aid gives a completely specified four-dimensional curvature-quadratic slice. Let \(M\) be any nonempty closed oriented smooth four-manifold, let \(a,b,c\in\R\), and define
\[
 S_g=(a|\operatorname{Rm}_g|^2+b|\operatorname{Ric}_g|^2+cR_g^2)\dd V_g.
\]
Here \(R_g\) is scalar curvature, \(W_g\) the Weyl tensor, and tensor norms use the full tensor contraction convention. Then the aid reports the equivalence
\[
 \left[\forall g\ \forall\varphi\in C^\infty(M),\quad
 \int_M S_{e^{2\varphi}g}=\int_M S_g\right]
 \quad\Longleftrightarrow\quad a+b+3c=0.
\]
The quantifier is every metric, not merely one fixed conformal class. Its witness lemma is reported to detect failure on every such manifold: if \(a+b+3c\ne0\), a suitable metric and infinitesimal conformal change give a nonzero first variation. No explicit witness construction is supplied in the summary.

For \(\Delta_g=\operatorname{div}_g\nabla_g\), fix
\[
 Q_4(g)=-\tfrac16\Delta_g R_g-\tfrac12|\operatorname{Ric}_g|^2+\tfrac16R_g^2,
 \qquad \beta=2a+b.
\]
With \(a+b+3c=0\), the reported pointwise density decomposition is
\[
 S_g=-2\beta Q_4(g)\dd V_g
       +a|W_g|^2\dd V_g
       -\tfrac\beta3(\Delta_g R_g)\dd V_g.
\]
The Weyl density is locally conformally invariant in dimension four. The last term is an exact divergence, since
\[
 (\Delta_g R_g)\dd V_g=\dd(\iota_{\nabla_gR_g}\dd V_g).
\]
Thus this slice has the stated \(Q+L+V\) form with its coefficients and Laplacian convention explicit. It concerns quadratic curvature expressions with no curvature derivatives in \(S_g\), although a divergence derivative appears in the decomposition.

The aid says the standard polynomial complete-contraction version of the general conjecture follows from Alexakis, but warns that the workshop's unqualified ``natural \(n\)-form'' may be broader. No precise theorem, allowed derivative orders, parity/dimensional hypotheses, naturality class, or full citation is supplied. The four-dimensional result must not silently be identified with a classification of all such potentially broader forms.
\subsection{Short English statement}
The excerpt states a conformal-invariant-density decomposition but omits the preceding definition of its input. The aid makes one four-dimensional slice precise: curvature-quadratic densities have conformally invariant integrals exactly when one linear coefficient relation holds, and then admit the displayed \(Q\)-curvature, Weyl, and divergence decomposition. The claimed general theorem covers a specified polynomial class rather than every possible interpretation of naturality.
\subsection{Sources}
\begin{itemize}
\item Local UnsolvedMath registry, permanent ID 20002012, exact source code catalogue entry; canonical \texttt{aim-geometry-notes.json}, \texttt{notes[349]}, tag \texttt{conjecture}; AIM Workshop Problem Lists, Geometry, workshop \emph{Conformal structure in geometry, analysis, and physics}, Conjecture 2. Created and updated 14 August 2026.
\item Supplied source URL \url{https://aimath.org/WWN/confstruct/confstruct.pdf}. The record gives only the short conjecture and does not reproduce ``as above''. The PDF was not accessed.
\item Alexakis is named for the standard polynomial complete-contraction version, without title, theorem number, publication entry, or separate URL. The broader natural-form scope requires source verification.
\item Attempt 1: \texttt{partial\_result}, problem status at run \texttt{partially\_solved}, rigor label \texttt{proved}, novelty confidence \texttt{low}; classification \texttt{PARTIAL-PROGRESS}, suggested difficulty L4, registry status \texttt{partially\_solved}. Candidate: coefficient-level four-dimensional classification and witness lemma. These are dated supplied assessments, not independent certification.
\item Ulam AI / UnsolvedMath Contributors; CC BY 4.0. Machine-generated research aid, not peer reviewed; independently verify before citation.
\end{itemize}
\end{document}
